\section{Relations, reductions, and proof objects}
\label{sec:encoding}

\begingroup\footnotesize

\subsection{How to read this section}

A paper-level encoded relation contains triples $(x,O;w)$: $x$ is the explicit
instance, $O$ is the implicit encoded oracle, and $w$ is the witness.  For a public
witness map $T$, the binding condition is
\[
 O=\Enc_{\mathcal C}(T(w)).
\]
Merkle compilation commits to $O$ with the public root
$\rho:=\Root(O)$; $\rho$ is not a proof field.
Every relation below is additionally parameterized by the fixed public
context $\chi:=(\mathsf{pp},t,s,W,g)$ and its derived dimensions.  The context
is omitted from displayed instances and is carried unchanged by every edge.

The edge notation is
\[
 \mathcal R_0\equiv_F\mathcal R_1
 \quad\text{deterministic retyping},\qquad
 \mathcal R_0\xRightarrow{\Pi}\mathcal R_1
 \quad\text{interactive reduction},
\]
\[
 \mathcal R_0\mathrel{\dashrightarrow}\mathcal R_1
 \quad\text{conditional extension or classification only}.
\]
Each subsection defines one source relation and lists its direct outgoing
reductions.  A proof fragment lists transmitted values only.  Challenges,
successor coefficients, exponentiated values, batching coefficients, and
work tables are verifier-derived.

\subsection{The application PIMSAT relations}

Let $S:=R_M[X]$, and let $M\subseteq S$ be a statement-bound finite free
$R_M$-submodule.  Fix a public instance $\bm x\in S^{\ell_x}$, a bounded
witness domain $\mathcal W_B\subseteq M^r$, ideals $I_a\lhd S$, and constraint
polynomials $p_j^{(a)}\in S[\bm X,\bm Y]$.  The paper first uses the
unencoded relation
\begin{align}
 \RPIMSAT^{\rm app}:=\Bigl\{(\bm x,\bot;\bm f):\;&
 \bm f\in\mathcal W_B,
 \quad p_j^{(a)}(\bm x,\bm f)\in I_a\quad(\forall a,j)\Bigr\}.
 \label{eq:rel-pimsat-unencoded}
\end{align}
Encoding the witness gives the encoded application relation
\begin{align}
 \RPIMSAT^{\rm enc}:=\Bigl\{(\bm x,O;\bm f):\;&
 \bm f\in\mathcal W_B,
 \quad O=\Enc_{\mathcal C}\!\left(
          \Pack(\pi_2(\bits(\bm f)))\right),\notag\\[-2pt]
 &p_j^{(a)}(\bm x,\bm f)\in I_a\quad(\forall a,j)\Bigr\}.
 \label{eq:rel-pimsat-app}
\end{align}

\paragraph{Direct reductions.}
\begin{itemize}

\item \textbf{Reduction 0: Compile the witness oracle.}
\[
 \RPIMSAT^{\rm app}
 \xRightarrow{\mathsf{Encode}}
 \RPIMSAT^{\rm enc}.
\]
The IOP supplies the encoded oracle $O$; Merkle compilation exposes
$\rho=\Root(O)$.  The constraint predicates are unchanged.

\item \textbf{Reduction 1: Reduce ideal membership to application PESAT.}
\[
 \RPIMSAT^{\rm enc}
 \xRightarrow{\Pi_{{\rm app},1}}
 \RPESAT^{\rm app}.
\]
\begin{itemize}
 \item \textbf{Proof:} $\pi_{{\rm app},1}$ carries the PIOP messages that turn
 ideal-membership constraints into field equations.
 \item \textbf{Result:} the verifier derives $K_{\rm app}$,
 $\psi_{\rm app}$, $\widehat{\bm x}$, and $\widehat{\bm p}$; $O$ is unchanged.
 \item \textbf{Scope:} application-supplied, before the F2Z PCS core.
\end{itemize}

\item \textbf{Classification: Equality and congruence constraints.}
$I_a=(0)$ gives equality; nonzero $I_a$ gives congruence or quotient-ring
constraints.  Both use the same reduction.

\end{itemize}

\subsection{The encoded application PESAT relation}

Fix the verifier-derived application field $K_{\rm app}$, a public map
$\psi_{\rm app}:M\to K_{\rm app}$, and projected constraint polynomials
$\widehat p_i\in K_{\rm app}[\bm X,\bm Y]$.  The encoded application PESAT
relation retains the original bounded witness:
\begin{align}
 \RPESAT^{\rm app}:=
 \Bigl\{((\widehat{\bm x},\widehat{\bm p}),O;\bm f):\;&
 \bm f\in\mathcal W_B,
 \quad O=\Enc_{\mathcal C}\!\left(
             \Pack(\pi_2(\bits(\bm f)))\right),\notag\\[-2pt]
 &\widehat p_i(\widehat{\bm x},\psi_{\rm app}(\bm f))=0
   \quad(\forall i)\Bigr\}.
 \label{eq:rel-pesat-app}
\end{align}

\paragraph{Direct reductions.}
\begin{itemize}

\item \textbf{Reduction 1: Linearize the application PESAT claim.}
\[
 \RPESAT^{\rm app}
 \xRightarrow{\Pi_{{\rm app},2}}
 \RLMFB.
\]
\begin{itemize}
 \item \textbf{Proof:} $\pi_{{\rm app},2}$ carries the application reduction
 from the projected constraints to one scalar linear claim.
 \item \textbf{Result:} the verifier derives $(\bm v,\mu)$ and $\psi$; the
 same witness and $O$ pass to $\RLMFB$.
 \item \textbf{Scope:} application-supplied.  Together,
 $\pi_{\rm PIOP}:=(\pi_{{\rm app},1},\pi_{{\rm app},2})$ is the complete
 application proof fragment.
\end{itemize}

\item \textbf{Classification: A single linear constraint is LMFB.}
If $\widehat{\bm p}$ is one linear polynomial, the relation is already
Lin-ModFieldBin; no extra message is needed.

\end{itemize}

\subsection{The Lin-ModFieldBin relation}

Let $(\gamma_j)_{j<D_M}$ be the fixed basis of $M$, and assume a
statement-bound bijection
$\bits:\mathcal W_B\leftrightarrow\bitsset^n$ with fixed reconstruction:
\[
 f_i=\sum_{k<B}\sum_{j<D_M}b_{i,k,j}2^k\gamma_j,
 \qquad b_{i,k,j}\in\bitsset,
 \qquad \bm b=\bits(\bm f).
\]
At this seam a higher-rank target has already been expanded into scalar
coordinates; write $\psi:M\to R$ for the resulting $\varphi$-semilinear map.
\Needspace{7\baselineskip}
The module relation is
\begin{align}
 \RLMFB:=\Bigl\{((\bm v,\mu),O;\bm f):\;&
 \bm f\in\mathcal W_B,
 \quad\bm v\in R^r,
 \quad\mu\in R,
 \quad O=\Enc_{\mathcal C}\!\left(
                 \Pack(\pi_2(\bits(\bm f)))\right),\notag\\[-2pt]
 &\sum_i v_i\psi(f_i)=\mu\quad\text{in }R\Bigr\}.
 \label{eq:rel-lmfb}
\end{align}

\paragraph{Direct reductions.}
\begin{itemize}

\item \textbf{Reduction 1: Deterministic Bitify.}

\[
 \RLMFB
 \equiv_{\mathsf{Bitify}}
 \RLBD,
 \qquad
 u_{i,k,j}:=
 v_i\,\varphi(\iota_{R_M}(2^k))\,\psi(\gamma_j).
\]
The witness is retyped as $\bm b:=\bits(\bm f)$; the same encoded oracle $O$
remains bound to it.  Bitify transmits nothing and performs no new Booleanity
test.

\item \textbf{Required preprocessing for a higher-rank target.}
If $N$ has rank greater than one, first expand it into scalar LMFB claims.
Then bitify each coordinate or use an application-defined batch; no canonical
batching proof is specified here.

\end{itemize}

\subsection{The Lin-BitsDual relation}

The core bit relation is
\begin{align}
 \RLBD:=\Bigl\{((\bm u,\mu),O;\bm b):\;&
 \bm b\in\bitsset^n,
 \quad\bm u\in R^n,
 \quad\mu\in R,
 \quad O=\Enc_{\mathcal C}(\Pack(\pi_2(\bm b))),\notag\\[-2pt]
 &\langle\iota_R(\bm b),\bm u\rangle_R=\mu\Bigr\}.
 \label{eq:rel-lbd}
\end{align}

Suppose the public coefficient has the paper's tensor form
$\bm u=\bm u^{(1)}\otimes\bm u^{(2)}$, and reshape the committed bits as
$\Bbit[c,i]$.  With canonical integer representatives
$\gamma_i$ of $u_i^{(1)}$, the prover exposes column values
\[
 \mu_c:=\sum_i\gamma_i\Bbit[c,i]\in\Z.
\]
The verifier immediately checks the tensor reconstruction
\begin{equation}
 \sum_c u_c^{(2)}\,\iota_R(\mu_c)=\mu.
 \label{eq:pi1-paper-reconstruction}
\end{equation}
It then derives $y_i:=g^{\gamma_i}$ and $\nu_c:=g^{\mu_c}$, obtaining
\[
 \prod_i\bigl(1+(y_i-1)\Bbit[c,i]\bigr)=\nu_c
 \qquad(\forall c).
\]

\paragraph{Direct reductions.}
\begin{itemize}

\item \textbf{Reduction 1: Produce the F2Z encoded PESAT claim.}
\[
 \RLBD
 \xRightarrow{\Pi_1}
 \RPESAT(\K,Q_{\Ftwo\!\Z},\mathcal C,\Pack).
\]
\begin{itemize}
 \item \textbf{Proof:} the paper sends $\pi_1=(\mu_c)_c$.  The verifier checks
 Equation~\eqref{eq:pi1-paper-reconstruction} before accepting the successor.
 \item \textbf{Result:} the verifier derives $\bm y=(g^{\gamma_i})_i$,
 $\bm\nu=(g^{\mu_c})_c$, and the PESAT instance; the bit witness and $O$
 are unchanged.
 \item \textbf{Scope:} the first core F2Z reduction.
\end{itemize}

\item \textbf{Profile restriction: weighted coefficient shape.}
The chunked mod-$q$ profile requires
\[
 u_{c,b,j}=\pi_q(w_b2^j)w'_c,
 \qquad
 0\le w_b<2^{q_{\rm bits}},\qquad w'_c\in\Fq.
\]
This restricts the public coefficient; it does not define another relation.

\end{itemize}

\subsubsection{Active-chunk instantiation of
  \texorpdfstring{$\Pi_1$}{Pi1}}

For the admitted weighted shape, write
\[
 w_b=\sum_{\ell=0}^{L-1}2^{c_w\ell}\omega_b^{(\ell)},
 \qquad
 u_c^{(\ell)}:=\sum_b\omega_b^{(\ell)}\INT(D(b,c)).
\]
The proof fragment is
\begin{equation}
 \pi_1^{\rm chunk}
 :=\bigl(\nu^{(\ell)},u^{(\ell)}\bigr)_{
       \ell\in\mathcal A\uparrow},
 \qquad
 \begin{aligned}
  u^{(\ell)}&=(u_c^{(\ell)})_{c\in\bitsset^s},\\
  \nu^{(\ell)}&=(g^{u_c^{(\ell)}})_{c\in\bitsset^s}.
 \end{aligned}
 \label{eq:pi1-chunk-proof}
\end{equation}
For each active chunk in increasing $\ell$, $\nu^{(\ell)}$ is a NARG record
and $u^{(\ell)}$ is a checked, non-absorbed hint.  Inactive chunks have the
derived value $u^{(\ell)}=0$.  The verifier rejects a missing, duplicate,
out-of-order, or noncanonical active pair; requires
$0\le u_c^{(\ell)}<2^{127}$; and checks
$g^{u_c^{(\ell)}}=\nu_c^{(\ell)}$ before its next squeeze.  Because
$\operatorname{ord}(g)=2^{128}-1$, the stated range makes the accepted exponent
unique.  The checked values reconstruct the paper's column messages:
\begin{equation}
 \mu_c
 :=\sum_{\ell=0}^{L-1}2^{c_w\ell}u_c^{(\ell)}
 =\sum_b w_b\INT(D(b,c)).
 \label{eq:chunk-reconstructs-mu}
\end{equation}
Inside $\Pi_1$, immediately after receiving the final active vector, the
verifier checks
\begin{equation}
 \sum_c w'_c\,\pi_q(\mu_c)=y.
 \label{eq:pi1-weighted-reconstruction}
\end{equation}
This is the tensor-reconstruction check of
Equation~\eqref{eq:pi1-paper-reconstruction}; it creates no named successor
relation and leaves no public scalar obligation for the opening phase.

For each active $\ell$, the verifier also derives
\[
i=(b\ll w)\mid j,\qquad
\omega_i^{(\ell)}=\omega_b^{(\ell)}2^j,\qquad
 y_i^{(\ell)}=g^{\omega_i^{(\ell)}}.
\]
Thus the single paper target is realized by the profile family
\begin{equation}
 \bigwedge_{\ell\in\mathcal A}
 \RPESAT(\K,Q_{\Ftwo\!\Z}^{(\ell)},\mathcal C,\Pack).
 \label{eq:chunk-pesat-family}
\end{equation}

\subsection{The encoded F2Z PESAT relation}

For one active profile chunk, define
\[
 Q_c^{(\ell)}(\Bbit)
 :=\prod_i\left(1+(y_i^{(\ell)}-1)\Bbit[c,i]\right)
   -\nu_c^{(\ell)}.
\]
The corresponding named paper relation is the encoded PESAT instance
\begin{align}
 &\RPESAT(\K,Q_{\Ftwo\!\Z}^{(\ell)},\mathcal C,\Pack)
 \notag\\[-2pt]
 &\quad:=
 \Bigl\{((\bm y^{(\ell)},\bm\nu^{(\ell)}),O;\Bbit):
 \Bbit\in\Ftwo^{2^s\times2^d},\ 
 O=\Enc_{\mathcal C}(\Pack(\Bbit)),\ 
 Q_c^{(\ell)}(\Bbit)=0\ (\forall c)\Bigr\}.
 \label{eq:rel-gp}
\end{align}
Omitting $\ell$ gives the paper's unchunked notation
$\RPESAT(\K,Q_{\Ftwo\!\Z},\mathcal C,\Pack)$.  The profile's ``grand
product'' is this PESAT relation, not a distinct formal relation.

\paragraph{Direct reductions and classifications.}
\begin{itemize}

\item \textbf{Reduction 1: Linearize with GKR.}
\[
 \RPESAT(\K,Q_{\Ftwo\!\Z},\mathcal C,\Pack)
 \xRightarrow{\Pi_{\rm red}}
 \RLin(\K,\mathcal C,\Pack).
\]
For the chunked family, $\Pi_{\rm red}$ is applied once per
$\ell\in\mathcal A$.
\begin{itemize}
 \item \textbf{Input:} the chunk's PESAT instance, the unchanged root
 $\rho=\Root(O)$, and the transcript after $\zeta^{(\ell)}$ and
 $e_0^{(\ell)}$; call this bundle $\mathcal I_{\rm GKR}^{(\ell)}$.
 \item \textbf{Proof:}
 \[
  (\tr',\mathit{pt}_\ell,\mu_\ell)
  \leftarrow
  \mathsf{GKR.Verify}
  (\mathsf{pp}_{\rm GKR},\mathcal I_{\rm GKR}^{(\ell)},
   \pi_{\rm red}^{(\ell)};\tr).
 \]
 $\pi_{\rm red}^{(\ell)}$ is one canonical GKR fragment.  Its profile fixes the
 messages, challenge order, encodings, tags, and rejection rules; the call
 consumes the whole fragment and returns $\tr'$.
 \item \textbf{Result:} the verifier derives $(\bm a_\ell,\mu_\ell)$ for
 $\RLin(\K,\mathcal C,\Pack)$, retaining the same $O$ and bit witness.  In the
 current profile, with
 $B_{\mathrm{bit},\K}:=\iota_{2\to\K}(\Bbit)$ coordinatewise,
 \[
  \bm a_\ell[z]=\eq(z,\mathit{pt}_\ell),
  \qquad
  \MLE{B_{\mathrm{bit},\K}}(\mathit{pt}_\ell)=\mu_\ell.
 \]
 This equation defines the LIN successor; it is not another paper relation.
 \item \textbf{Scope:} this edge is solid once the GKR profile fixes a
 byte-exact transcript and proves the reduction.  Its error is
 $\varepsilon_{\rm red}^{(\ell)}$ per active chunk.
\end{itemize}

\item \textbf{Classification: Degree-zero, zero-ideal PIMSAT.}
With $S=\K[X]$, $S_{\le0}=\K$, and $I=(0)$, membership means exactly
$Q_c^{(\ell)}=0$.  Boolean witnesses keep their bit domain; unrestricted
$\K$ variables also require $Z^2-Z=0$.

\end{itemize}

\subsection{The encoded LIN relation}

Let $\iota_{2\to\K}:\Ftwo\hookrightarrow\K$ act coordinatewise and put
$\bm b_{\K}:=\iota_{2\to\K}(\pi_2(\bm b))$.  The paper's linear target is
\begin{align}
 \RLin(\K,\mathcal C,\Pack):=
 \Bigl\{((\bm a,\mu),O;\bm b):\;&
 \bm b\in\bitsset^n,
 \quad\bm a\in\K^n,
 \quad\mu\in\K,\notag\\[-2pt]
 &O=\Enc_{\mathcal C}(\Pack(\pi_2(\bm b))),
 \quad\langle\bm a,\bm b_{\K}\rangle_{\K}=\mu\Bigr\}.
 \label{eq:rel-lin-pack}
\end{align}

\paragraph{Direct reductions.}
\begin{itemize}

\item \textbf{Reduction 1: Run the linear-relation IOP.}
\[
 \RLin(\K,\mathcal C,\Pack)
 \xRightarrow{\Pi_{\rm iopp}}
 \mathcal R_{\rm triv}.
\]
\begin{itemize}
 \item \textbf{Proof:} $\pi_{\rm iopp}$ proves the derived $(\bm a,\mu)$ claim
 against the public root $\rho=\Root(O)$.
 \item \textbf{Result:} acceptance or rejection; no public relation remains.
 \item \textbf{Scope:} the final solid paper-level edge.  The profile below
 gives one concrete implementation.
\end{itemize}

\end{itemize}

\subsubsection{Current profile's internal expansion of
  \texorpdfstring{$\Pi_{\rm iopp}$}{Pi-iopp}}

The chunked recursive-opening profile implements $\Pi_{\rm iopp}$ on the
MLE-shaped LIN family emitted by $\Pi_{\rm red}$ as follows.  These are
internal claim transformations, not additional solid edges in the paper's
relation graph.
\begin{enumerate}

\item \textbf{Per-claim ring switch.}
For each active $\ell$, the prover transmits
\[
 s^{(\ell)}=(s_v^{(\ell)})_{v<128},
 \qquad
 s_v^{(\ell)}:=\MLE{B_{\mathrm{bit},v}}(r_{\hi}^{(\ell)}).
\]
After all such vectors are absorbed, the verifier squeezes one shared $r''$
and derives a packed claim
\begin{equation}
 O=\Enc_{\mathcal C}(P),
 \qquad
 \langle B_\ell,P\rangle_{\K}=\beta_\ell.
 \label{eq:profile-packed-claim}
\end{equation}
The witness is $P=\Pack(\pi_2(\bm b))$ and $\rho=\Root(O)$ is unchanged.  This step
has error $\varepsilon_{\rm ring}^{(\ell)}$ on the admitted MLE-shaped source;
it does not specify a general ring switch for arbitrary LIN coefficients.

\item \textbf{Batch.}
The verifier squeezes fresh
$\bm\eta\in\K^{|\mathcal A|}$ and derives, without a proof field,
\[
 B:=\sum_{\ell\in\mathcal A}\eta_\ell B_\ell,
 \qquad
 \beta:=\sum_{\ell\in\mathcal A}\eta_\ell\beta_\ell.
\]
The resulting internal claim is
$O=\Enc_{\mathcal C}(P)$ and
$\langle B,P\rangle_{\K}=\beta$, with batching error
$\varepsilon_{\rm batch}$.

\item \textbf{Recursive opening.}
The prover transmits $\pi_{\rm open}=\pi_{\rm lig}$, one profile-defined
recursive constrained-code proof containing its internal proximity messages
and authenticated openings.  The verifier derives all opening challenges and
accepts only if the proof verifies against $(\rho,B,\beta)$.  This completes
$\Pi_{\rm iopp}$ with error $\varepsilon_{\rm open}$.

\end{enumerate}

The intended paper ring-switch construction is not used as an independent
solid edge here: the three operations above are the selected profile's
internal implementation of the already named
$\RLin(\K,\mathcal C,\Pack)\xRightarrow{\Pi_{\rm iopp}}\mathcal R_{\rm triv}$
edge.

\subsection{The virtual Lin-BitsDual extension}

For a public $\Ftwo$-linear map $T:\Ftwo^n\to\Ftwo^{n'}$ and
$h_0\in\Ftwo^{n'}$, let
$\operatorname{lift}_2:\Ftwo\to\{0,1\}\subset\Z$ act coordinatewise and put
$\bar h=h_0+T(\pi_2(\bm b))$.  Define
\begin{align}
 \RLBD^{\virt}:=\Bigl\{((T,\bm a_q,\mu,h_0),O;\bm b):\;&
 \bm b\in\bitsset^n,
 \quad\bm a_q\in\Fq^{n'},
 \quad\mu\in\Fq,\notag\\[-2pt]
 &O=\Enc_{\mathcal C}(\Pack(\pi_2(\bm b))),
 \quad
 \langle\pi_q(\operatorname{lift}_2(\bar h)),\bm a_q\rangle_{\Fq}
 =\mu\Bigr\}.
 \label{eq:rel-lbd-virtual}
\end{align}

\paragraph{Direct reductions.}
\begin{itemize}

\item \textbf{Conditional virtual reduction.}
\[
 \RLBD^{\virt}
 \overset{\Pi_{\rm partial\text{-}core}}{\dashrightarrow}
 \mathcal R_{h\text{-}{\rm Lin}}^{\virt}.
\]
\begin{itemize}
 \item \textbf{Proof:} $\pi_{\rm virtual}$ carries the chosen partial-core and
 GKR messages; the virtual table is not transmitted.
 \item \textbf{Result:} the verifier derives $(\bm a_{\K},\mu')$.
 \item \textbf{Scope:} available only after a virtual profile defines the
 proof format, transcript, verifier, and error bound.
\end{itemize}

\end{itemize}

Let $h_{0,\K}:=\iota_{2\to\K}(h_0)$ and let $T_{\K}$ be the scalar extension
of $T$.  Only $\bm b$ is committed.  The virtual linear seam is
\begin{align}
 \mathcal R_{h\text{-}{\rm Lin}}^{\virt}:=
 \Bigl\{((T,h_0,\bm a_{\K},\mu'),O;\bm b):\;&
 \bm b\in\bitsset^n,
 \quad\bm a_{\K}\in\K^{n'},
 \quad\mu'\in\K,\notag\\[-2pt]
 &O=\Enc_{\mathcal C}(\Pack(\pi_2(\bm b))),\notag\\[-2pt]
 &\langle\bm a_{\K},h_{0,\K}+T_{\K}(\bm b_{\K})\rangle_{\K}
 =\mu'\Bigr\}.
 \label{eq:rel-virtual-lin}
\end{align}
Its public adjoint rewrite transmits nothing:
\[
 \mathcal R_{h\text{-}{\rm Lin}}^{\virt}
 \equiv_{\mathsf{Adj}_{T_{\K}^{*}}}
 \RLin(\K,\mathcal C,\Pack),
\]
\[
 \bm a_b:=T_{\K}^{*}(\bm a_{\K}),
 \qquad
 \mu_b:=\mu'-\langle\bm a_{\K},h_{0,\K}\rangle_{\K}.
\]
This exact rewrite need not preserve the MLE coefficient shape admitted by the
current internal ring switch.  A virtual profile must therefore supply a
compatible $\Pi_{\rm iopp}$ adapter or opening; the direct profile rejects the
virtual source.

\subsection{The complete relation graph and proof object}

Let $\mathcal R_{\rm triv}:=\{(\bot,\bot;\bot)\}$.  The paper-level relation
graph is
\begin{equation}
\begin{aligned}
 \RPIMSAT^{\rm app}
 &\xRightarrow{\mathsf{Encode}}
 \RPIMSAT^{\rm enc},
 \\[1pt]
 \RPIMSAT^{\rm enc}
 &\xRightarrow{\Pi_{{\rm app},1}}
 \RPESAT^{\rm app}
 \xRightarrow{\Pi_{{\rm app},2}}
 \RLMFB,
 \\[1pt]
 \RLMFB
 &\equiv_{\mathsf{Bitify}}
 \RLBD
 \xRightarrow{\Pi_1}
 \RPESAT(\K,Q_{\Ftwo\!\Z},\mathcal C,\Pack)
 \\[1pt]
 \RPESAT(\K,Q_{\Ftwo\!\Z},\mathcal C,\Pack)
 &\xRightarrow{\Pi_{\rm red}}
 \RLin(\K,\mathcal C,\Pack)
 \xRightarrow{\Pi_{\rm iopp}}
 \mathcal R_{\rm triv}.
\end{aligned}
\label{eq:solid-reduction-graph}
\end{equation}
$\mathsf{Encode}$ and Bitify add no proof fields; chunking and recursive opening
refine the displayed edges.  The semantic proof object is
\begin{equation}
 \pi_{\mathrm{F2Z}}
 =\left(
   \pi_{\rm PIOP},
   \bigl(\nu^{(\ell)},u^{(\ell)},\pi_{\rm red}^{(\ell)}\bigr)_{
         \ell\in\mathcal A\text{ increasing}},
   (s^{(\ell)})_{\ell\in\mathcal A\text{ increasing}},
   \pi_{\rm lig}
 \right).
 \label{eq:proof-object}
\end{equation}
Its replay schedule is specified in Section~\ref{sec:interaction}.  Public
inputs are external to the proof; other omitted values are derived.  The v1
byte profile partitions transmitted fields, in replay order, as
\begin{equation}
 \pi_{\rm wire}:=(N_{\rm F2Z},H_{\rm F2Z}),\qquad
 \begin{aligned}
  \NargLen(\pi)&:=|N_{\rm F2Z}|,&
  \HintLen(\pi)&:=|H_{\rm F2Z}|,\\
  \PayloadLen(\pi)&:=|N_{\rm F2Z}|+|H_{\rm F2Z}|,&
    \WireLen(\pi)&:=
    |\HostEncode(N_{\rm F2Z},H_{\rm F2Z})|.
 \end{aligned}
 \label{eq:proof-channel-lengths}
\end{equation}
Benchmarks \must{} report all four lengths.  Section~\ref{sec:wire-container}
fixes their byte meaning, the outer 72-byte container, and the exact record
counts and exhaustion rules.

With edge errors $\varepsilon_{{\rm app},1}$, $\varepsilon_{{\rm app},2}$,
$\varepsilon_{\rm red}^{(\ell)}$, $\varepsilon_{\rm ring}^{(\ell)}$,
$\varepsilon_{\rm batch}$, and $\varepsilon_{\rm open}$,
\[
 \varepsilon_{\rm F2Z}
 \le \varepsilon_{{\rm app},1}+\varepsilon_{{\rm app},2}
    +\sum_{\ell\in\mathcal A}
       \bigl(\varepsilon_{\rm red}^{(\ell)}
             +\varepsilon_{\rm ring}^{(\ell)}\bigr)
    +\varepsilon_{\rm batch}+\varepsilon_{\rm open}.
\]

\endgroup
